\section{Introduction}

Stock-price forecasting remains a common applied machine-learning problem. A typical study gathers daily OHLCV data, adds technical indicators, trains a neural network, and reports RMSE, MAE, or MAPE on a held-out period. More recent versions of this workflow replace recurrent networks with transformer-based or foundation-style time-series models. However, a central evaluation problem remains: absolute stock prices are highly persistent. If the target is the future price level, the no-training rule that the future close equals the latest observed close is often extremely hard to beat.

This observation changes the order of evidence required for a credible paper. Before asking whether one architecture is better than another, the first question is whether any learned model adds value beyond latest-close persistence on the same chronological test labels. If the answer is no, then claims about feature engineering, architectural novelty, or model sophistication are not supported for this target.

The original version of this project was a conventional applied stock-prediction study using LSTM/GRU models and technical indicators. Expanded audit experiments invalidated that positive framing. The current paper therefore reframes the work as a baseline-first diagnostic audit:

\begin{quote}
Do selected supervised and zero-shot time-series models reliably beat latest-close persistence in absolute stock-price level forecasting?
\end{quote}

Under the audited protocol, the answer is largely negative. Supervised neural models almost never clear the persistence gate. Chronos-T5-small zero-shot forecasting produces more local wins, especially at longer horizons, but remains slightly worse than persistence in median RMSE ratio. This is a stronger and more honest contribution than claiming a fragile positive result from a narrow model comparison.

The paper makes three contributions:
\begin{itemize}
\item It defines a cross-market baseline-first protocol for absolute stock-price forecasting, using matched chronological evaluation and a mandatory persistence gate.
\item It audits recurrent, linear, transformer-style, multiscale, and zero-shot foundation-model forecasters across 165 dataset-series memberships, two price modes, and three horizons.
\item It reports a negative result that is useful for applied researchers: architecture comparisons should be interpreted only after a model beats no-training persistence.
\end{itemize}
